% Vietnamese translation for OI Public Library
% Source-File: IOI中国国家候选队论文/2006/王赟/王赟.ppt
\documentclass[11pt]{article}
\usepackage{oipl-vn}

\newcommand{\Oh}{\mathcal{O}}

\title{Đồ thị Trie\\{\large Ghi chú theo slide}}
\author{Wang Yun}
\date{2006}

\begin{document}
\maketitle

\origtitle{Trie graph construction, use, and improvement}
\sourcepath{IOI China candidate team papers / 2006 / Wang Yun / Wang Yun.ppt}

\section{Phạm vi}

Bộ trình chiếu đi kèm bài viết về cây trie, đồ thị trie và các cải tiến bộ
nhớ. Ghi chú này dựa trên bản bài viết đã dịch.

\section{Từ trie tới tự động}

Cây trie lưu một tập từ theo tiền tố. Khi thêm liên kết hậu tố và các cạnh
chuyển trạng thái còn thiếu, ta thu được đồ thị trie, tức tự động Aho--Corasick
cho bài toán khớp nhiều mẫu. Mỗi trạng thái biểu diễn một tiền tố đã đọc, còn
liên kết hậu tố đưa ta tới hậu tố dài nhất vẫn là tiền tố của một mẫu.

\section{Từ cấm}

Trong bài dò từ cấm, mỗi nút có thể được đánh dấu nguy hiểm nếu chuỗi từ gốc
tới nút kết thúc bằng một mẫu cấm. Dấu nguy hiểm được truyền qua liên kết hậu
tố. Khi quét văn bản, chỉ cần đi theo đồ thị; gặp trạng thái nguy hiểm nghĩa
là văn bản chứa một mẫu cấm.

\section{Dùng thông tin trên đồ thị}

Đồ thị trie còn dùng được cho đếm số lần xuất hiện, tìm vị trí mẫu trong ma
trận chữ, hoặc xây đồ thị an toàn gồm các trạng thái không nguy hiểm. Sau khi
có đồ thị an toàn, nhiều bài trở thành kiểm tra đường đi, chu trình hoặc quy
hoạch động trên trạng thái.

\section{Cải tiến bộ nhớ}

Lưu bảng chuyển đầy đủ tốn \(\Oh(L a)\), với \(L\) là tổng độ dài mẫu và \(a\)
là kích thước bảng chữ. Nếu \(a\) lớn hoặc trie thưa, có thể lưu cạnh hiện có
bằng danh sách, băm, hoặc nén các chuyển mặc định qua liên kết hậu tố để giảm
bộ nhớ.

\end{document}
